\begin{problem}[第35届全国中学生物理竞赛复赛][引力场运动与轨道]{88}
一、假设地球是一个质量分布各向同性的球体。从地球上空离地面高度为 $h$ 的空间站发射一个小物体，该物体相对于地球以某一初速度运动，初速度方向与其到地心的连线垂直。已知地球半径为 $R$，质量为 $M$，引力常量为 $G$。地球自转及地球大气的影响可忽略。
\begin{subproblem}
若该物体能绕地球做周期运动, 其初速度的大小应满足什么条件?
\end{subproblem}
\begin{subproblem}
若该物体的初速度大小为 $v_{0}$ ，且能落到地面，求其落地时速度的大小和方向（即速度与其水平分量之间的夹角）、以及它从开始发射直至落地所需的时间。

已知：对于 $c < 0$ ， $\Delta = b^{2} - 4ac > 0$ ，有
\begin{equation}
\int \frac {x \mathrm{d} x}{\sqrt {a + b x + c x ^ {2}}} = \frac {\sqrt {a + b x + c x ^ {2}}}{c} - \frac {b}{2 (- c) ^ {3 / 2}} \arcsin \frac {2 c x + b}{\sqrt {\Delta}} + C
\label{eq:1.p88}
\end{equation}
式中 $C$ 为积分常数。
\end{subproblem}
\end{problem}
